\documentclass[10pt]{article}
\usepackage[T1]{fontenc}
\usepackage{lmodern,amsmath,amssymb,amsthm,geometry}
\usepackage[hidelinks]{hyperref}
\geometry{margin=1in}
\newtheorem{theorem}{Theorem}
\newtheorem{lemma}[theorem]{Lemma}
\newcommand{\Bit}{\operatorname{Bit}}
\newcommand{\I}{\mathsf I}
\newcommand{\M}{\mathsf M}
\newcommand{\U}{\mathsf U}
\newcommand{\E}{\mathsf E}
\newcommand{\N}{\mathsf N}
\title{The two membership heads are independently necessary\\for their own derivability in the PA196 calculus}
\author{Research note}
\date{13 November 2026}
\begin{document}
\maketitle

\begin{abstract}
A tempting state reduction deletes one of PA196's two membership
rewriting heads and tries to recover it using the other and pointed
separation.  Two explicit interpretations show that this particular
strategy cannot work: each head fails in an interpretation validating
all the remaining exact native rules.  The interpretation is sensitive
to the numerical labels of variables, which is permitted when proving
independence of this deliberately restricted numeric calculus.  We also
give a nontrivial consequence of one-sided rewriting, namely ordinary
same-variable universal elimination, to distinguish the obstruction
from the failure of the simplest specialization arguments.
\end{abstract}

\section{The question}

The exact definitions are those of
\href{https://github.com/jamesyc/BusyBeaver/blob/407cdc356f6b586635e71ac92be394255f4bfea4/PA196/Logic/Basic.lean}{\texttt{PA196/Logic/Basic.lean}}.
Up to exchange of antecedents, the two heads are
\begin{align*}
 L(a,b,d)&:\quad \E(a,b)\to(\M(a,d)\to\M(b,d)),\\
 R(a,b,d)&:\quad \E(a,b)\to(\M(d,a)\to\M(d,b)).
\end{align*}
They are the paper's head 3 and head 4, respectively.  The remaining
rules are modus ponens, closed generalization, heads 1, 2, 5 and 6,
the restricted propositional head $\rho_8$, and the entire guarded gate,
including equality vacuity and membership vacuity.

An independence interpretation need only validate these exact numeric
rules.  It need not validate all unpostulated rules of ordinary
first-order logic, especially uniform substitution of variable labels
inside membership atoms.  Assuming that additional invariance would
assume precisely the phenomenon to be derived.

\section{A two-point permutation}

Let $\sigma:\mathbb N\to\mathbb N$ interchange $2$ and $3$ and fix
all other numbers.  It is an involution.  Both identities
\[
 \neg\Bit(\sigma(t),t),\qquad
 \neg\Bit(t,\sigma(t))
\]
hold for every $t$.  Away from $2,3$ they are the usual absence of the
$t$th bit in $t$.  At $2,3$ they follow by direct calculation:
$\Bit(3,2)=\Bit(2,3)=0$.
Set $f_1=g_1=\sigma$ and $f_i=g_i=\mathrm{id}$ for $i\ne1$.

\section{Head 3 is not derivable without itself}

Given an assignment $\nu:\mathbb N\to\mathbb N$, interpret atoms by
\[
 \M(a,b)^\nu:=\Bit\bigl(f_a(\nu(a)),\nu(b)\bigr),\qquad
 \E(a,b)^\nu:=(\nu(a)=\nu(b)).
\]
Interpret implication classically, and interpret $\U(a,B)$ as
universal quantification of the $a$th coordinate of the assignment.
Use the exact numeric decoder, including code $0$ as truth and the
reversed quantifier pairing.  Since $f_0=\mathrm{id}$, code
$1=\M(0,0)$ is false under every assignment.  Hence all the derived
connectives have their intended classical meanings, even when
quantification binds coordinate $0$.

The quantifier heads 1, 2 and 6, the equality witness head 5, and
$\rho_8$ are valid for ordinary reasons.  Head 4 is valid: if
$\nu(a)=\nu(b)$, the two membership formulas in $R(a,b,d)$ test
the same digit of the same number $f_d(\nu(d))$.  Modus ponens and
closed generalization preserve validity.

It remains to check the entire gate, because arbitrary separation
is the possible source of a hidden derivation.  Let $c\ne a,b$ and
let $\varphi$ be the decoded property with original code $c$.
The source-label bound shows that changing coordinate $c$ does not
change $\varphi$.  Write
\[
 P(y):=\varphi^{\nu[b:=y]}.
\]
If $a\ne b$, take the finite set
\[
 Z=\{y:(\Bit(f_a(\nu(a)),y)\vee y=\nu(b))\wedge P(y)\}.
\]
Let $m=\sum_{y\in Z}2^y$, and choose $s=f_c^{-1}(m)$.  Then
$\M(c,b)$ under $\nu[c:=s,b:=y]$ is precisely $y\in Z$.  This
establishes the separation conjunct and the distinguished-point
clause.  If $a=b$, the first antecedent of the separation conjunct
is $\Bit(f_b(y),y)$, which is always false.  In this case use just
$Z=\{\nu(b):P(\nu(b))\}$; the remaining subset condition and the
point clause still hold.  Equality and membership vacuity are valid
because $c\ne a,b$.  Thus the whole bundled gate is valid.

But head 3 fails.  Use the actual instance
\[
 \M(1,2)\to\bigl(\E(1,0)\to\M(0,2)\bigr)
\]
and the assignment $\nu(1)=\nu(0)=2$, $\nu(2)=0$.  Its first
antecedent is $\Bit(3,0)$, which is true; its equality antecedent is
true; its conclusion is $\Bit(2,0)$, which is false.

\begin{theorem}
Head 3 is not derivable from the remaining exact PA196 native rules.
This remains true if primitive equality reflexivity, symmetry and
transitivity are added, since equality in the counterinterpretation
is actual equality of natural numbers.
\end{theorem}
\begin{proof}
Every remaining rule preserves validity in the interpretation, and
the displayed instance is invalid.
\end{proof}

\section{Head 4 is not derivable without itself}

Use instead
\[
 \M(a,b)^\nu:=\Bit\bigl(\nu(a),g_b(\nu(b))\bigr),\qquad
 \E(a,b)^\nu:=(\nu(a)=\nu(b)).
\]
The same argument verifies the logical heads and rules.  Head 3 is
now valid, because it replaces the container value by an equal
value and leaves the digit index $g_d(\nu(d))$ unchanged.

For the gate with $a\ne b$, consider the finite set of indices
\[
 Z=\{g_b(y):(\Bit(\nu(a),g_b(y))\vee y=\nu(b))\wedge
                    \varphi^{\nu[b:=y]}\}.
\]
It is finite because $g_b$ is a bijection and a natural number has
only finitely many set bits.  Take the witness
$s=\sum_{z\in Z}2^z$ for coordinate $c$.  Injectivity of $g_b$
ensures that no point with the wrong property enters the selected
set.  If $a=b$, the first antecedent is $\Bit(y,g_b(y))$, always
false; the singleton-point construction suffices again.
All other parts of the gate are as before.

Head 4 fails at
\[
 \E(0,1)\to\bigl(\M(2,0)\to\M(2,1)\bigr)
\]
with $\nu(0)=\nu(1)=2$, $\nu(2)=4$.  The membership antecedent
is $\Bit(4,2)$ and is true; the conclusion is $\Bit(4,3)$ and is
false.

\begin{theorem}
Head 4 is not derivable from the remaining exact PA196 native rules,
even with primitive equality reflexivity, symmetry and transitivity.
\end{theorem}

\section{A useful positive consequence of just head 4}

This section uses head 4, head 5, the other logical heads and the gate,
but never head 3.  Vacuity for arbitrary formulas at a fresh label
follows as in the upstream construction; that proof uses the gate,
negation, implication and quantifier commutation, not either
membership rewrite head.

\begin{lemma}[Negative-membership specialization]
For $r\ne x$ and $z\ne x$,
\[
 \U(x,\N\M(r,x))\to\N\M(r,z)
\]
is derivable using head 4.
\end{lemma}
\begin{proof}
The equality witness is $\exists x_x\E(z,x)$.  Head 4 gives
$\E(z,x)\to(\M(r,z)\to\M(r,x))$.  Under the assumption
$\M(r,z)$, its vacuity at $x$ and universal distribution give
$\U(x,\E(z,x)\to\M(r,x))$.  Combining this with the universal
negative premise gives $\U(x,\N\E(z,x))$, contradicting the
witness.  Discharge the assumption.
\end{proof}

\begin{lemma}[Same-variable universal elimination]
For every label $x$ and formula code $A$,
\[
 \U(x,A)\to A
\]
is derivable without head 3.
\end{lemma}
\begin{proof}
Let $c$ be the code of $\N A$.  If $c=x$, use the code of
$\N\N\N A$ instead; it is different and is propositionally
equivalent to $\N A$.  The gate at $(x,x,c)$ gives a witness
$x_c$ with a point clause implying $\M(c,x)$ from $\N A$ and
a separation clause implying
$\U(x,\M(c,x)\to\N A)$.
Assume $\U(x,A)$ and $\N A$.  The clauses give $\M(c,x)$ and
$\U(x,\N\M(c,x))$.

Choose a label $z$ fresh for all displayed formulas and distinct from
$0,x,c$.  Negative-membership specialization gives $\N\M(c,z)$.
Head 4 gives $\E(x,z)\to(\M(c,x)\to\M(c,z))$.  Thus under the
present assumptions $\N\E(x,z)$ follows.  Generalize at $z$;
the assumptions are vacuous there.  This contradicts head 5's
$\exists x_z\E(x,z)$.  Finally eliminate the separation witness
at $c$, which is fresh in the desired conclusion, and discharge
$\N A$ by classical logic.
\end{proof}

This result does not cross from a bound label $x$ to a different
free label $y$.  Such a crossing is exactly where an attempted
proof of the missing membership head needs an additional,
unjustified substitution or alpha-conversion argument.  The two
counterinterpretations show that no completion of that route can
produce the missing head from the unchanged remaining rules.

\section{Limits of the conclusion}

The countermodels exclude a specific reduction justification:
``delete one head because the other heads already derive it.''
They do not prove that a machine searching the weaker calculus
cannot still be equiconsistent with PA through a different
interpretation.  Nor do they establish any lower bound on Turing
machine states.  A new PA interpretation of a one-head calculus
would be a separate mathematical construction.

\end{document}
